\originalsection{第3次习题课}
\sourceinfo{MIT 18.04, Spring 2018, Recitation 3; 题面 \nolinkurl{mit18_04s18_recit3-handout.pdf}, PDF第1页; 解答 \nolinkurl{mit18_04s18_recit3-solutions.pdf}, PDF第1--3页. 课程作者 Jeremy Orloff, 页内署名 Vishesh Jain, 2018, CC BY-NC-SA 4.0}
\begin{exercise}
\textbf{原题 Rec3.1。}\label{ass:rec3-1}
设 $f:\C\to\C$ 解析, 写 $f(x+\ii y)=u(x,y)+\ii v(x,y)$. 假设 $u,v\in C^2$, 即二阶及以下偏导存在且连续. 证明 $f'$ 也解析.
\end{exercise}
\begin{solution}
据官方解答翻译补全 (解答PDF第1页). 由 Cauchy--Riemann 方程 $u_x=v_y$, $u_y=-v_x$, 有 $f'=u_x+\ii v_x$. 令 $U=u_x$, $V=v_x$. 二阶偏导连续允许交换混合偏导, 因而 $U_x=u_{xx}=v_{yx}=v_{xy}=V_y$, $U_y=u_{xy}=u_{yx}=-v_{xx}=-V_x$. $U,V$ 的一阶偏导连续并满足 Cauchy--Riemann 方程, 故 $f'$ 在每一点复可微, 即解析. 此处按原题的 $C^2$ 假设完成论证; 实际上解析函数自动有任意阶导数.
\end{solution}
\begin{exercise}
\textbf{原题 Rec3.2。}\label{ass:rec3-2}
2.1. 证明 $\int\bar z\dd z$ 在 $\C$ 中不与路径无关, 为什么不违背复线积分基本定理? 2.2. 对每个 $n\in\Z$, 计算以原点为圆心的单位圆 $\gamma$ 上的 $\int_\gamma z^n\dd z$, 检查与基本定理的相容性. 2.3. 若 $\gamma$ 换成所围闭圆盘不含原点的圆, 答案是否改变? 圆均取正向.
\end{exercise}
\begin{solution}
据官方解答翻译补全 (解答PDF第1页). 2.1沿 $z=\ee^{\ii t}$, $0\le t\le2\pi$, 得 $\int_\gamma\bar z\dd z=\int_0^{2\pi}\ee^{-\ii t}\ii\ee^{\ii t}\dd t=2\pi\ii\ne0$. 与常值闭路径相比积分不同. 基本定理要求被积函数具有复原函数; 若 $F'=\bar z$, 对所有闭路的积分都应为0, 这里正说明这样的原函数不存在.

2.2直接参数化得 $\int_\gamma z^n\dd z=\ii\int_0^{2\pi}\ee^{\ii(n+1)t}\dd t$, 所以 $n=-1$ 时为 $2\pi\ii$, 其余整数时为0. 当 $n\ne-1$ 时, $z^{n+1}/(n+1)$ 在被积函数的定义域上是单值原函数, 因而基本定理给出0. 当 $n=-1$ 时, $1/z$ 在穿孔平面没有单值原函数, 多值对数不能直接用于闭路基本定理. 2.3所有积分均为0: 被积函数在所围闭圆盘的某邻域解析, 可以用 Cauchy 定理; 对 $1/z$ 也可在该圆盘邻域取解析对数, 再应用基本定理.
\end{solution}
\begin{exercise}
\textbf{原题 Rec3.3。}\label{ass:rec3-3}
回忆 $\cos(x+\ii y)=\cos x\cosh y-\ii\sin x\sinh y$. 3.1. 在 $\mathcal R=\{x+\ii y:0<x<\pi\}$ 中, 水平线和竖直线的像是什么? 限制在 $\mathcal R$ 上的余弦映射是否一一? 3.2. 向 $\mathcal R$ 添加半线 $x=0,y\ge0$ 及 $x=\pi,y>0$, 得到 $\mathcal R_1$. 求其像, 并判断是否仍然一一. 3.3. $\mathcal R_1$ 给出多值反余弦的一种取值选择, 此分支在反余弦定义域的割线是什么?
\end{exercise}
\begin{solution}
据官方解答翻译补全 (解答PDF第1--3页); 退化情形、单射与满射论证以及3.2--3.3的修正为编者补充 (AI). 写 $w=u+\ii v$. 固定 $x=x_0\ne\pi/2$, 得 $u^2/\cos^2x_0-v^2/\sin^2x_0=1$, 且 $u$ 与 $\cos x_0$ 同号, 所以像为双曲线的一个完整分支. 当 $x_0=\pi/2$ 时公式分母为0, 实际像是整个虚轴 $w=-\ii\sinh y$. 固定 $y=y_0\ne0$, 得 $u^2/\cosh^2y_0+v^2/\sinh^2y_0=1$, 但 $0<x<\pi$ 仅给出半条椭圆, $v$ 的符号与 $y_0$ 相反, 两端点不包含. 当 $y_0=0$ 时退化为实区间 $(-1,1)$.

若 $\cos z=\cos\zeta$, 用指数式或余弦差公式得 $z=\zeta+2k\pi$ 或 $z=-\zeta+2k\pi$. 两个实部均在 $(0,\pi)$ 时, 前一种仅允许 $k=0$, 后一种不可能, 所以在 $\mathcal R$ 上单射. 为求像, 令 $q=\ee^{\ii z}$, 方程为 $q^2-2wq+1=0$. 当 $w$ 不在实轴外侧两条射线上时, 两根非实且互为倒数, 恰有一根的辐角在 $(0,\pi)$; 若 $w\in(-1,1)$, 两根在单位圆上, 仍恰有一根位于上半圆. 对这根取 $z=\arg q-\ii\log|q|\in\mathcal R$, 得到
\[
\cos\mathcal R=\C\setminus\bigl(( -\infty,-1]\cup[1,\infty)\bigr).
\]
这里非实 $w$ 不会有实根, 互倒的两根虚部异号, 保证上述选择唯一. 当 $w\ge1$ 或 $w\le-1$ 时, 两根为实根, 因而没有 $z\in\mathcal R$.

在新添半线上, $\cos(\ii y)=\cosh y\in[1,\infty)$, $\cos(\pi+\ii y)=-\cosh y\in(-\infty,-1)$. 两半线的像与内部像互不相交, 且各自一一, 所以 $\mathcal R_1$ 仍为一一映射, 像恰为 $\C\setminus\{-1\}$. \textbf{官方解答称像为整个复平面, 但原题第二条半线要求 $y>0$, 因而遗漏 $z=\pi$, 也就遗漏 $w=-1$.} 若把该条件改为 $y\ge0$, 才得到全平面的一一取值选择.

3.3的解析反函数定义在开割域 $\C\setminus(( -\infty,-1]\cup[1,\infty))$, 割线为两条含端点的实射线. 原解答写两条不含端点的射线; 作为描述跳跃位置可以如此说, 但 $\pm1$ 是分支点, 解析分支的定义域还须排除它们. $\mathcal R_1$ 在割线上选定一侧的值, 并不使反函数在割线两侧连续. 例如 $w>1$ 两侧极限分别为 $\pm\ii\operatorname{arcosh}w$, 所以发生跳跃. $w<-1$ 两侧相应为 $\pi\pm\ii\operatorname{arcosh}(-w)$.
\begin{center}
\begin{tikzpicture}[scale=.65,>=Stealth]
\begin{scope}[xshift=-4cm]
\draw[->](-.3,0)--(3.6,0) node[right]{$x$};\draw[->](0,-2.4)--(0,2.4)node[above]{$y$};
\foreach \x in {.35,.7,1.05,1.57,2.09,2.44,2.79}{\draw[blue!70] (\x,-1.6)--(\x,1.6);}
\foreach \y in {-1.5,-1,-.5,0,.5,1,1.5}{\draw[orange!80!black] (.04,\y)--(3.1,\y);}
\draw[dashed](3.1416,-2.2)--(3.1416,2.2);\node[below] at(3.1416,0){$\pi$};\node at(1.57,-2.9){$z$ 平面};
\end{scope}
\begin{scope}[xshift=3cm]
\draw[->](-3.5,0)--(3.5,0)node[right]{$u$};\draw[->](0,-2.8)--(0,2.8)node[above]{$v$};
\foreach \a in {.35,.7,1.05,1.57,2.09,2.44,2.79}{\draw[blue!70,smooth,domain=-1.6:1.6,samples=65]plot({cos(deg(\a))*cosh(\x)},{-sin(deg(\a))*sinh(\x)});}
\foreach \b in {-.5,-1,-1.5,.5,1,1.5}{\draw[orange!80!black,smooth,domain=.04:3.1,samples=65]plot({cos(deg(\x))*cosh(\b)},{-sin(deg(\x))*sinh(\b)});}
\draw[very thick](-.999,0)--(.999,0);\draw[fill=white](-1,0)circle(1.5pt) node[below]{$-1$};\draw[fill=white](1,0)circle(1.5pt)node[below]{$1$};\node at(0,-3.6){$w=\cos z$};
\end{scope}
\end{tikzpicture}
\end{center}

\end{solution}
